\documentclass[11pt]{article}
\usepackage[margin=1in]{geometry}
\usepackage{enumitem}
% =============================================================================
% Preambles/header.tex — shared LaTeX/Beamer preamble for this project
%
% Use from a Beamer lecture by prepending:
%     \input{header}
% and compiling with TEXINPUTS=../Preambles:$TEXINPUTS (the /compile-latex
% skill does this automatically).
%
% PALETTE CONTRACT. Color names defined here MUST match the names in
% Quarto/theme-template.scss so Beamer and Quarto renderings use the same
% palette. The scripts/check-palette-sync.sh script enforces this.
%
% To customize: edit the HEX values below AND the matching $variable in
% Quarto/theme-template.scss, then run scripts/check-palette-sync.sh.
% =============================================================================

% -----------------------------------------------------------------------------
% Engine detection + encoding
%
% Our /compile-latex skill uses XeLaTeX, where inputenc is unnecessary and
% can produce encoding warnings. Only load inputenc under pdfTeX.
% -----------------------------------------------------------------------------
\usepackage{iftex}
\ifPDFTeX
  \usepackage[utf8]{inputenc}
\fi

\usepackage{amsmath, amssymb, amsthm}
\usepackage{booktabs}
\usepackage{graphicx}

% -----------------------------------------------------------------------------
% PALETTE — mirrors Quarto/theme-template.scss variable names.
% Load xcolor and define colors BEFORE anything that references them
% (notably hyperref, hypersetup, and the Beamer color assignments).
% -----------------------------------------------------------------------------
\usepackage{xcolor}

\definecolor{primary-blue}{HTML}{012169}      % headings, accents
\definecolor{primary-gold}{HTML}{B9975B}      % emphasis, borders
\definecolor{highlight-yellow}{HTML}{F2A900}  % markers, alerts
\definecolor{light-bg}{HTML}{E8EDF5}          % title / section backgrounds
\definecolor{jet}{HTML}{1A1A1A}               % body text

% Semantic colors (used in diagrams and annotations)
\definecolor{positive}{HTML}{15803D}          % good / observed / identified
\definecolor{negative}{HTML}{B91C1C}          % bad / problematic / bias
\definecolor{neutral}{HTML}{525252}           % reference / context

% Highlight accents (match .hi-* classes in the SCSS)
\definecolor{hi-slate}{HTML}{314F4F}
\definecolor{hi-green}{HTML}{2E8B57}
\definecolor{hi-red}{HTML}{C41E3A}

% Semantic aliases so skill templates and snippets can write
% \textcolor{accent}{...} without hard-coding "primary-blue".
\colorlet{accent}{primary-blue}
\colorlet{accent2}{primary-gold}

% -----------------------------------------------------------------------------
% Hyperref — loaded AFTER xcolor + palette so link colors resolve.
% -----------------------------------------------------------------------------
\usepackage{hyperref}
\hypersetup{colorlinks=true, linkcolor=primary-blue, urlcolor=primary-blue,
            citecolor=primary-blue, pdfborder={0 0 0}}

% -----------------------------------------------------------------------------
% Course code — set once per deck (after \input{header}, before \begin{document})
% via \coursecode{CS401}. Shown in the footer on every slide so a deck's course
% is identifiable without opening the title page. Empty by default.
% -----------------------------------------------------------------------------
\makeatletter
\newcommand{\coursecode}[1]{\gdef\@coursecodetext{#1}}
\gdef\@coursecodetext{}
\makeatother

% -----------------------------------------------------------------------------
% Beamer theme assignments (only apply under Beamer).
%
% We detect Beamer via \@ifclassloaded{beamer}, which is the documented,
% non-internal way. This must be wrapped in \makeatletter/\makeatother
% because \@ifclassloaded contains an @-catcode character.
% -----------------------------------------------------------------------------
\makeatletter
\@ifclassloaded{beamer}{%
  \usefonttheme{professionalfonts}
  \setbeamercolor{structure}{fg=primary-blue}
  \setbeamercolor{frametitle}{fg=primary-blue}
  \setbeamercolor{title}{fg=primary-blue}
  \setbeamercolor{subtitle}{fg=primary-gold}
  \setbeamercolor{author}{fg=primary-blue}
  \setbeamercolor{institute}{fg=neutral}
  \setbeamercolor{date}{fg=primary-gold}
  \setbeamercolor{itemize item}{fg=primary-blue}
  \setbeamercolor{itemize subitem}{fg=primary-gold}
  \setbeamercolor{alerted text}{fg=primary-gold}
  \setbeamercolor{block title}{fg=white, bg=primary-blue}
  \setbeamercolor{block body}{bg=light-bg}
  % Semantic block variants: block=definition/concept (blue), exampleblock=
  % worked example (green/positive), alertblock=key takeaway or caution (gold).
  % Keep to <=2 colored boxes per slide across ALL three variants (INV-7).
  \setbeamercolor{block title example}{fg=white, bg=positive}
  \setbeamercolor{block body example}{bg=light-bg}
  \setbeamercolor{block title alerted}{fg=white, bg=primary-gold}
  \setbeamercolor{block body alerted}{bg=light-bg}

  % Minimal footer: course code (left) + page X / N (right), no navigation clutter
  \setbeamertemplate{navigation symbols}{}
  \setbeamertemplate{footline}{%
    \color{neutral}\footnotesize\hspace{0.7em}\@coursecodetext\hfill
    \insertframenumber/\inserttotalframenumber\hspace{0.7em}\vspace{0.3em}}
}{}
\makeatother

% -----------------------------------------------------------------------------
% TikZ setup — libraries the snippet gallery uses
% -----------------------------------------------------------------------------
\usepackage{tikz}
\usetikzlibrary{arrows.meta, positioning, calc, decorations.pathreplacing,
                fit, shapes.geometric, backgrounds}

% Shared TikZ styles — snippets and hand-written diagrams can pick these up
% via `\begin{tikzpicture}[every node/.style={...}]` or by naming specific
% styles (e.g., `\node[dag-node] (X) at (0,0) {X};`).
\tikzset{
  % Node styles for causal DAGs and similar
  dag-node/.style = {
    draw=primary-blue, fill=light-bg, rounded corners=2pt,
    minimum width=1.2cm, minimum height=0.9cm, align=center, font=\small
  },
  decision-node/.style = {
    draw=primary-gold, fill=white, diamond, aspect=1.8,
    minimum width=1.8cm, minimum height=1.2cm, align=center, font=\small,
    inner sep=1pt
  },
  % Edge styles — observed vs counterfactual
  observed-edge/.style = {-{Latex[length=2.5mm]}, thick, draw=jet},
  counterfactual-edge/.style = {-{Latex[length=2.5mm]}, thick, draw=neutral, dashed},
  confound-edge/.style = {-{Latex[length=2.5mm]}, thick, draw=negative, dashed},
  % Data-point styles for plots
  observed-dot/.style = {fill=primary-blue, draw=primary-blue, circle, inner sep=1.2pt},
  counterfactual-dot/.style = {fill=white, draw=neutral, circle, inner sep=1.2pt}
}

% -----------------------------------------------------------------------------
% Convenience macros
% -----------------------------------------------------------------------------
\newcommand{\muted}[1]{\textcolor{neutral}{#1}}
\newcommand{\key}[1]{\textcolor{primary-gold}{\textbf{#1}}}
\newcommand{\good}[1]{\textcolor{positive}{#1}}
\newcommand{\bad}[1]{\textcolor{negative}{#1}}

% Standout frame macro: full-bleed dark background for section transitions.
% Beamer-only (uses \begin{frame}/\setbeamercolor/\usebeamerfont) -- do not
% call from an article-class file (e.g. Notes/<CODE>/*-notes.tex). \muted,
% \key, \good, \bad, and \coursecode above are plain xcolor/gdef and are
% safe to use from either class; verified when Notes/article-class support
% was added.
% Expand: \transitionslide{Your transition title}
\newcommand{\transitionslide}[1]{%
  \begingroup
  \setbeamercolor{background canvas}{bg=primary-blue}
  \begin{frame}[plain]
    \centering
    \vfill
    {\usebeamerfont{title}\color{white}\Large #1\par}
    \vfill
  \end{frame}
  \endgroup
}

% Section divider frame (Beamer-only), an alternative to \transitionslide with a
% darker two-line style: near-black (jet) background, a small white label line
% above a larger gold title, and a thin gold rule along the bottom edge.
% Expand: \sectiondivider{Section 2}{Pointers}
\newcommand{\sectiondivider}[2]{%
  \begingroup
  \setbeamercolor{background canvas}{bg=jet}
  \begin{frame}[plain]
    \vspace*{3.0cm}
    {\color{white}\ttfamily\large #1\par}
    \vspace{0.5cm}
    {\color{primary-gold}\LARGE #2\par}
    \begin{tikzpicture}[remember picture, overlay]
      \draw[primary-gold, line width=2.5pt]
        ($(current page.south west)+(0,0.1)$) -- ($(current page.south east)+(0,0.1)$);
    \end{tikzpicture}
  \end{frame}
  \endgroup
}

% -----------------------------------------------------------------------------
% Channel watermark — "Concepts That Click" (YouTube).
%
% Article-class files (CompetitiveExam Q/A sets, Notes, Assignments) get a
% light diagonal draft watermark on every page plus a centered footer naming
% the channel. Beamer decks get a subtle TikZ background watermark instead
% (the footline already carries the course code). Centralized here so every
% MCQ PDF is branded without per-file edits.
% -----------------------------------------------------------------------------
\makeatletter
\@ifclassloaded{article}{%
  \usepackage{draftwatermark}
  \SetWatermarkText{Concepts That Click}
  \SetWatermarkScale{1}
  \SetWatermarkFontSize{2cm}
  \SetWatermarkLightness{0.86}
  \SetWatermarkAngle{45}
  \usepackage{fancyhdr}
  \pagestyle{fancy}
  \fancyhf{}
  \fancyfoot[C]{\footnotesize\textcolor{neutral}{Concepts That Click~|~YouTube: \href{https://www.youtube.com/@concepts.thatclick}{@concepts.thatclick}}}
  \renewcommand{\headrulewidth}{0pt}
  \renewcommand{\footrulewidth}{0pt}
  \fancypagestyle{plain}{%
    \fancyhf{}%
    \fancyfoot[C]{\footnotesize\textcolor{neutral}{Concepts That Click~|~YouTube: \href{https://www.youtube.com/@concepts.thatclick}{@concepts.thatclick}}}%
    \renewcommand{\headrulewidth}{0pt}%
    \renewcommand{\footrulewidth}{0pt}%
  }
}{}
\@ifclassloaded{beamer}{%
  \setbeamertemplate{background}{%
    \begin{tikzpicture}[remember picture, overlay]
      \node[rotate=45, opacity=0.055, align=center]
        at (current page.center)
        {\Huge\textcolor{neutral}{Concepts That Click}};
    \end{tikzpicture}%
  }
}{}
\makeatother 
\coursecode{JKSSB}
\title{Lesson 32 Practice: PowerPoint Animation and Transition}
\author{JKSSB Computer Awareness}
\date{\today}
\begin{document}
\maketitle
\noindent\textbf{Provenance:} Q3 is a real previous-year item
(\textit{[PYQ WO 2026 Q80, key:provisional]}); Q17 is adapted from the JA 2026
SmartArt item (\textit{[PYQ-adapted]}); the remaining items are original
JKSSB-pattern questions (\textit{[Original]}).

\begin{enumerate}[label=\textbf{Q\arabic*.}, leftmargin=*]
\item \textit{[Original]} In PowerPoint, an animation is applied to:
  \begin{enumerate}[label=\Alph*.]
    \item the whole presentation at once
    \item an object or text on a slide
    \item the change from one slide to the next
    \item the file name of the presentation
  \end{enumerate}
\item \textit{[Original]} In PowerPoint, a transition is applied to:
  \begin{enumerate}[label=\Alph*.]
    \item a single object on a slide
    \item the movement from one slide to the next
    \item the notes pane
    \item the spelling and grammar checker
  \end{enumerate}
\item \textit{[PYQ WO 2026 Q80, key:provisional]} Which key starts the slide
  show \textbf{from the beginning}?
  \begin{enumerate}[label=\Alph*.]
    \item F5
    \item F7
    \item F9
    \item Shift+F5
  \end{enumerate}
\item \textit{[Original]} Which shortcut starts the slide show \textbf{from the
  current slide}?
  \begin{enumerate}[label=\Alph*.]
    \item F5
    \item Ctrl+F5
    \item Shift+F5
    \item Alt+F5
  \end{enumerate}
\item \textit{[Original]} In PowerPoint, \textbf{Alt+F5} is commonly used to:
  \begin{enumerate}[label=\Alph*.]
    \item end the slide show
    \item start the presentation in Presenter View
    \item spell-check the current slide
    \item open the Notes pane
  \end{enumerate}
\item \textit{[Original]} Which key \textbf{ends} a running slide show?
  \begin{enumerate}[label=\Alph*.]
    \item F5
    \item Esc
    \item Enter
    \item Tab
  \end{enumerate}
\item \textit{[Original]} Which of the following is \textbf{NOT} an animation
  effect category in PowerPoint?
  \begin{enumerate}[label=\Alph*.]
    \item Entrance
    \item Exit
    \item Transition
    \item Emphasis
  \end{enumerate}
\item \textit{[Original]} Which animation effect moves an object \textbf{along a
  defined path} on the slide?
  \begin{enumerate}[label=\Alph*.]
    \item Entrance
    \item Exit
    \item Emphasis
    \item Motion Path
  \end{enumerate}
\item \textit{[Original]} Which animation effect is used to make an object
  \textbf{appear} on the slide?
  \begin{enumerate}[label=\Alph*.]
    \item Entrance
    \item Exit
    \item Emphasis
    \item Motion Path
  \end{enumerate}
\item \textit{[Original]} An effect that draws attention to an object already
  visible on the slide, \textbf{without removing it}, is called:
  \begin{enumerate}[label=\Alph*.]
    \item Entrance
    \item Exit
    \item Emphasis
    \item Transition
  \end{enumerate}
\item \textit{[Original]} Which transition smoothly moves and animates matching
  objects as one slide changes into the next?
  \begin{enumerate}[label=\Alph*.]
    \item Fade
    \item Morph
    \item Push
    \item Wipe
  \end{enumerate}
\item \textit{[Original]} The Notes pane in PowerPoint is mainly used to:
  \begin{enumerate}[label=\Alph*.]
    \item print several slides on one page
    \item store speaker notes for the presenter
    \item animate objects on the slide
    \item change the theme colours
  \end{enumerate}
\item \textit{[Original]} A handout in PowerPoint is intended to:
  \begin{enumerate}[label=\Alph*.]
    \item control the slide transitions
    \item give the audience printed slides, often several per page
    \item start the slide show from the first slide
    \item edit the Slide Master
  \end{enumerate}
\item \textit{[Original]} Which view displays small thumbnails of all slides so
  they can be reordered?
  \begin{enumerate}[label=\Alph*.]
    \item Slide Sorter view
    \item Reading View
    \item Notes Page view
    \item Slide Show view
  \end{enumerate}
\item \textit{[Original]} Which of the following is \textbf{NOT} a standard view
  in PowerPoint?
  \begin{enumerate}[label=\Alph*.]
    \item Normal view
    \item Slide Sorter view
    \item Reading View
    \item Formula View
  \end{enumerate}
\item \textit{[Original]} Evaluate the statements:
  \begin{enumerate}[label=\Roman*.]
    \item An animation is applied to objects on a slide.
    \item A transition is applied between slides.
    \item Shift+F5 starts the slide show from the current slide.
  \end{enumerate}
  \begin{enumerate}[label=\Alph*.]
    \item I and II only
    \item II and III only
    \item I and III only
    \item I, II, and III
  \end{enumerate}
\item \textit{[PYQ-adapted]} When a SmartArt graphic is converted to text, what
  is lost and must be re-applied or re-entered?
  \begin{enumerate}[label=\Alph*.]
    \item nothing; the design is preserved exactly
    \item the graphical layout --- it becomes a plain bulleted list
    \item the file name of the presentation
    \item the slide transitions
  \end{enumerate}
\item \textit{[Original]} During a running slide show, pressing the
\textbf{B} key will:
  \begin{enumerate}[label=\Alph*.]
    \item end the slide show
    \item turn the screen black
    \item go back to the previous slide
    \item open Presenter View
  \end{enumerate}
\item \textit{[Original]} Which key is correctly matched with its PowerPoint
  function?
  \begin{enumerate}[label=\Alph*.]
    \item F5 --- spelling and grammar check
    \item F7 --- start the slide show
    \item Esc --- end the slide show
    \item Alt+F5 --- end the slide show
  \end{enumerate}
\item \textit{[Original]} Which pairing of a PowerPoint feature with its function
  is correct?
  \begin{enumerate}[label=\Alph*.]
    \item Morph --- a transition between similar slides
    \item Notes pane --- an animation effect
    \item Motion Path --- a transition between slides
    \item Handout --- starts the show from the current slide
  \end{enumerate}
\end{enumerate}
\end{document}
